% Part: normal-modal-logic
% Chapter: filtrations
% Section: euclidean-filtrations

\documentclass[../../../include/open-logic-section]{subfiles}

\begin{document}

\olfileid{nml}{fil}{euc}

\olsection{యూక్లిడియన్ నమూనాల వడపోతలు}

\olref[acc]{sec}లోని విధానం యూక్లిడియన్, లేదా
సీరియల్-యూక్లిడియన్ నమూనాల విషయంలో పనిచేయదు.
\olref{fig:ser-eucl} పైభాగంలోని నమూనాను చూడండి;
అది యూక్లిడియన్ కూడా, సీరియల్ కూడా. $\Gamma =
\{p, \Box p \}$ అనుకుందాం. $\Gamma$ ద్వారా వడపోత
చేసినప్పుడు $[w_1] = [w_3]$: ఎందుకంటే~$\Gamma$పై
ఏకీభవించే లోకాలు $w_1$, $w_3$ మాత్రమే.
\olref{fig:ser-eucl2}లో చూపినట్లుగా, ఏ వడపోతలోనైనా
$\mModel{M}$ నుంచి సంక్రమించిన బాణం కూడా ఉంటుంది.
ఆ నమూనా యూక్లిడియన్ కాదు. దానిని యూక్లిడియన్‌గా
మార్చడానికి అదనపు బాణాలు చేర్చలేం. $[w_2]$,
$[w_4]$ మధ్య రెండు దిశల్లో బాణాలు, ఆపై $[w_2]$,
$[w_5]$ మధ్య కూడా రెండు దిశల్లో బాణాలు చేర్చాల్సి
వస్తుంది. కానీ~$w_2$ వద్ద $\Box
p$ సత్యం కావాలి;~$w_5$ వద్ద $p$ అసత్యం.
\sourcecorrection{OLTENMLFILEUC-004}{స్థిర మూలం
  రెండో జత బాణాల చివరలను $w_2$, $w_5$గా
  రాసింది. ఇక్కడ బాణాలు వడపోత నమూనాలోనివి
  కనుక చివరలు $[w_2]$, $[w_5]$ కావాలి;
  తరువాతి సత్య పరీక్ష మాత్రం అసలు లోకాల
  $w_2$, $w_5$ వద్దనే ఉంటుంది.}

\begin{figure}[htpb]
  \centering
  \begin{tikzpicture}[modal]
    \node[world] (w1)
          [label={left: \mFalse{p}},
            label={below: $\mSat{{}}{\Box p}$}] {$w_1$};
    \node[world] (w2)
          [label={right: \mTrue{p}},
            label = {below: $\mSat{{}}{\Box p}$},
            right=of w1] {$w_2$};
    \draw[->] (w1) to (w2);
    \draw[reflexive above] (w2) to (w2);
    \node[world] (w3)
          [label={left: \mFalse{p}},
            label={below: $\mSat{{}}{\Box p}$},
            below=of w1] {$w_3$};
    \node[world] (w4)
          [label={right:\mTrue{p}},
            label={below: $\mSat/{{}}{\Box p}$},
            right=of w3] {$w_4$};
    \draw[->] (w3) to (w4);
    \draw[reflexive above] (w4) to (w4);
    \node[world] (w5)
          [label={right: \mFalse{p}},
            label=below:$\mSat/{{}}{\Box p}$,
            right=of w4] {$w_5$};
    \draw[->, bend left] (w4) to (w5);
    \draw[->, bend left] (w5) to (w4);
    \draw[reflexive above] (w5) to (w5);
  \end{tikzpicture}
  \caption{సీరియల్, యూక్లిడియన్ నమూనా.}
  \ollabel{fig:ser-eucl}
\end{figure}

\begin{figure}[ht]
  \centering
  \begin{tikzpicture}[modal,every text node part/.style={align=center}]
    \node[world] (w1)
          [label={left: \mFalse{p}},
            label={right:$[w_1]=[w_3]$},
            label={below: $\mSat{{}}{\Box p}$}] {$[w_1]$};
    \node[world] (w2)
         [label={right: \mTrue{p}},
           label = {below: $\mSat{{}}{\Box p}$},
           right=of w1, yshift=1.5cm] {$[w_2]$};
    \draw[->] (w1) to (w2);
    \draw[reflexive above] (w2) to (w2);
    \node[world] (w4)
         [label={right:\mTrue{p}},
           label={below: $\mSat/{{}}{\Box p}$},
           right=of w1,yshift=-1.5cm] {$[w_4]$};
    \draw[->] (w1) to (w4);
    \draw[reflexive above] (w4) to (w4);
    \node[world] (w5)
         [label={right: \mFalse{p}},
             label=below:$\mSat/{{}}{\Box p}$,
             right=of w4] {$[w_5]$};
    \draw[->, bend left] (w4) to (w5);
    \draw[->, bend left] (w5) to (w4);
    \draw[reflexive above] (w5) to (w5);
  \end{tikzpicture}
  \caption{\olref{fig:ser-eucl}లోని నమూనా వడపోత.}
  \ollabel{fig:ser-eucl2}
\end{figure}
\sourcecorrection{OLTENMLFILEUC-001}{స్థిర మూలంలోని
  మొదటి చిత్రంలో $w_2$కు బయటకు వెళ్లే బాణం లేదు;
  కాబట్టి అది ప్రకటించినట్లు సీరియల్ కాదు,
  $w_1Rw_2$ ఉన్నందున యూక్లిడియన్ కూడా కాదు.
  అవసరమైన $w_2$ స్వబాణాన్ని మొదటి చిత్రంలో,
  దాని వర్గానికి సంక్రమించే స్వబాణాన్ని రెండో
  చిత్రంలో చేర్చాం. $p$, $\Box p$ గుర్తులు మారవు.}

ముఖ్యంగా, యూక్లిడియన్ వడపోత పొందడానికి,
\tetoken{ఉపసూత్రాల}{!!{formula}s} పరంగా సంవృతమైన
ఏకపక్ష $\Gamma$ సమితుల ద్వారా వడపోతలను
పరిగణించడం సరిపోదు. బదులుగా
\emph{మోడల్ సంయోజకాల పరంగా సంవృతమైన}
$\Gamma$ సమితులను తీసుకోవాలి
(\olref[pre]{defn:modallyclosed} చూడండి).
అలాంటి శూన్యం కాని వాక్య సమితులు అనంతం;
అందువల్ల వాటి వల్ల సంబంధిత వ్యవస్థకు పరిమిత
నమూనా ధర్మం లేదా నిర్ణేయత వెంటనే లభించవు.
\sourcecorrection{OLTENMLFILEUC-002}{స్థిర మూలం
  మోడల్ సంయోజకాల పరంగా సంవృతమైన ప్రతి సమితీ
  అనంతమని చెప్పింది. ఖాళీ సమితి కూడా ఆ
  నిర్వచనాన్ని తృప్తిపరుస్తుంది; అది పరిమితమే.
  అందుకే ``శూన్యం కాని'' అని స్పష్టం చేశాం.}

\begin{thm}\ollabel{thm:modal-closed-filt}
  $\Gamma$ మోడల్ సంయోజకాల పరంగా సంవృతం,
  $\mModel{M}=\tuple{W,R,V}$ ఒక నమూనా,
  $\mModel{M^*} = \tuple{W^*,R^*,V^*}$ అనేది
  $\mModel{M}$కు చేసే అత్యంత స్థూల వడపోత అనుకుందాం.
  \begin{enumerate}
  \item $\mModel{M}$ సౌష్ఠవమైతే
    $\mModel{M^*}$ కూడా సౌష్ఠవమే.
  \item $\mModel{M}$ సంక్రామకమైతే
    $\mModel{M^*}$ కూడా సంక్రామకమే.
  \item $\mModel{M}$ యూక్లిడియన్ అయితే
    $\mModel{M^*}$ కూడా యూక్లిడియన్.
  \end{enumerate}
\end{thm}

\begin{proof}
\sourcecorrection{OLTENMLFILEUC-003}{స్థిర మూల
  సిద్ధాంతంలో అంశాల క్రమం సౌష్ఠవం, సంక్రామకత్వం,
  యూక్లిడియన్; కానీ నిరూపణలో ముందుగా సంక్రామకత్వం,
  తరువాత యూక్లిడియన్, చివర సౌష్ఠవం ఉన్నాయి.
  అదే నిరూపణను, అదే రెండు వ్యాయామాలను సిద్ధాంత
  అంశాల క్రమానికి అనుగుణంగా ఇక్కడ అమర్చాం.}
\begin{enumerate}
  \item వ్యాయామం. అన్ని సౌష్ఠవ నమూనాల్లో
    \Ax{B}, $\Ax{B_\Diamond}$ చెల్లుబాటవుతాయని
    వాడండి.
  \item $\mModel{M^*}$ అత్యంత స్థూల వడపోత అయితే,
    నిర్వచనం ప్రకారం $R^*[u][v]$ అయితే,
    అప్పుడు మాత్రమే $C_1(u,v)$. సంక్రామకత్వం కోసం
    $C_1(u,v)$, $C_1(v,w)$ అనుకుందాం; $C_1(u,w)$
    చూపాలి. \iftag{prvBox}{$\mSat{M}{\Box !A}[u]$
      అనుకుందాం. సంక్రామక నమూనాలన్నింటిలో
      \Ax{4} చెల్లుబాటవుతుంది కనుక
      $\mSat{M}{\Box\Box!A}[u]$. మోడల్ సంవృతత
      వల్ల $\Box\Box!A \in \Gamma$ కూడా; కనుక
      $C_1(u,v)$ వల్ల $\mSat{M}{\Box!A}[v]$;
      $C_1(v,w)$ వల్ల $\mSat{M}{!A}[w]$ కూడా.}{}\iftag{prvDiamond}{$\mSat{M}{!A}[w]$
      అనుకుందాం. మోడల్ సంవృతత వల్ల
      $\Diamond !A \in \Gamma$ కనుక $C_1(v,w)$
      నుంచి $\mSat{M}{\Diamond !A}[v]$.
      మోడల్ సంవృతత వల్ల
      $\Diamond\Diamond!A \in \Gamma$ కూడా;
      కనుక $C_1(u,v)$ నుంచి
      $\mSat{M}{\Diamond\Diamond !A}[u]$.
      సంక్రామక నమూనాలన్నింటిలో
      $\Ax{4}_\Diamond$ చెల్లుబాటవుతుంది కనుక
      $\mSat{M}{\Diamond!A}[u]$.}{}
  \item వ్యాయామం. అన్ని యూక్లిడియన్ నమూనాల్లో
    \Ax{5}, $\Ax{5_\Diamond}$ చెల్లుబాటవుతాయని
    వాడండి.
\end{enumerate}
\end{proof}

\begin{prob}
  \olref[nml][fil][euc]{thm:modal-closed-filt}
  నిరూపణను పూర్తి చేయండి.
\end{prob}

\end{document}
